% AUTO-GENERATED by corpus/make_appendix_table.py — do not edit by hand.
\begingroup
\renewcommand{\thetable}{A.\arabic{table}}
\setcounter{table}{0}
\begin{landscape}
\scriptsize
\begin{longtable}{@{}l p{2.6cm} p{3.4cm} p{1.8cm} p{5.0cm} p{2.1cm} p{6.2cm}@{}}
\caption{Included studies: characteristics and design (excerpt: 28 of 67 studies --- the 25 in the highest-ranked journals (AJG 4*, 4, 3) plus the first three AJG 2 studies; within a rating by year). Source: author's compilation.}\label{tab:coding-a1}\\
\toprule
ID & Study & Journal (AJG) & Country & Sample & Method & AI investment measure \\
\midrule
\endfirsthead
\multicolumn{7}{@{}l}{\tablename~\thetable{} (continued)}\\
\toprule
ID & Study & Journal (AJG) & Country & Sample & Method & AI investment measure \\
\midrule
\endhead
\bottomrule
\endfoot
S09 & Mishra et al. (2022) & Journal of the Academy of Marketing Science (4*) & USA & 19,000 firm-year observations, US-listed firms (COMPUSTAT + EDGAR 10-Ks), simultaneous equations; PSM robustness (1,801 treatment / 2,251 control firms) & panel econometrics & 10-K text (share of AI keywords = AI focus) \\
\addlinespace[2pt]
S14 & Krakowski et al. (2023) & Strategic Management Journal (4*) & multi-country & same chess players across conventional, centaur and engine tournaments (controlled competitive setting) & panel econometrics & AI (engine) adoption in competitive interaction \\
\addlinespace[2pt]
S17 & Babina et al. (2024) & Journal of Financial Economics (4*) & USA & 1,993 US Compustat firms; resume-based AI workforce measure (Cognism), robust to job-postings measure; IV = university AI graduate supply & panel econometrics & resume-based (share of AI-skilled workers as firm AI investment) \\
\addlinespace[2pt]
S07 & Leoni et al. (2022) & International Journal of Operations and Production Management (4) & Italy & 120 senior executives of manufacturing firms, PLS-SEM & survey-SEM & survey construct (AI adoption/use) \\
\addlinespace[2pt]
S45 & Adams et al. (2026) & Journal of Monetary Economics (4) & USA & US firms matched to Lightcast job postings 2010Q1-2024Q1 (40,000+ job boards); AI-pricing jobs identified & panel econometrics & job postings (AI-skill + pricing keyword = AI pricing adoption) \\
\addlinespace[2pt]
S46 & Agag et al. (2026) & Tourism Management (4) & UK & 1,206 UK listed tourism/travel/hospitality firms, 7,539 firm-years 2018-2024; two-step system GMM & panel econometrics & AI human capital (AI-skilled workforce measure) \\
\addlinespace[2pt]
S52 & Filieri et al. (2026) & Journal of Product Innovation Management (4) & China & Chinese listed manufacturers + WTO tariff records + customs data 2015-2022; IV analysis & panel econometrics & archival AI integration measure (tariff-induced) \\
\addlinespace[2pt]
S02 & Baabdullah et al. (2021) & Industrial Marketing Management (3) & Saudi Arabia & 392 B2B SMEs, cross-section, CB-SEM & survey-SEM & survey construct (acceptance of AI practices) \\
\addlinespace[2pt]
S03 & Chatterjee et al. (2021) & Industrial Marketing Management (3) & India & 349 managers from 16 BSE-listed firms (11 manufacturing, 5 service), response rate 51.6\%, PLS-SEM & survey-SEM & survey construct (AI-CRM implementation) \\
\addlinespace[2pt]
S04 & Chatterjee et al. (2022) & Journal of Business Research (3) & India & 312 usable manager responses from 14 BSE-listed firms (10 large/4 SME; telecom, IT, automobile, retail, pharma), Jan-Feb 2020, RR 47.3\%, PLS-SEM & survey-SEM & survey construct (adoption of AI-embedded CRM system) \\
\addlinespace[2pt]
S05 & Hossain et al. (2022) & Industrial Marketing Management (3) & Bangladesh & 257 usable manager responses (1,953 approached, 278 qualified), RMG export manufacturers, research-firm panel/Qualtrics, random sampling; time-lagged robustness subsample n=93 & survey-SEM & survey construct (AI adoption on top of marketing analytics platform) \\
\addlinespace[2pt]
S06 & Lee et al. (2022) & Technovation (3) & South Korea & 160 high-tech ventures (from 1,248 ministry list, 300 responses), survey + audited financials & panel econometrics & survey (AI-adoption intensity) linked to financial records \\
\addlinespace[2pt]
S08 & Lui et al. (2022) & Annals of Operations Research (3) & multi-market & 119 AI investment announcements by 62 listed firms, Factiva search Jan 2015 - Feb 2019 & event study & announcements (AI investment announcements) \\
\addlinespace[2pt]
S12 & Czarnitzki et al. (2023) & Journal of Economic Behavior and Organization (3) & Germany & 5,851 firms (409 AI users, \textasciitilde{}7\%), German CIS 2018 (ZEW); cross-section OLS/2SLS-IV + panel subset with IV-FE, entropy balancing & panel econometrics & survey (CIS AI adoption dummy + intensity of AI use in business processes) \\
\addlinespace[2pt]
S15 & Samadhiya et al. (2023) & Industrial Marketing Management (3) & India & in-depth interviews with senior B2B managers + PLS-SEM on 142 responses & mixed & survey construct (AI-based partner relationship management implementation) \\
\addlinespace[2pt]
S16 & Wu \& Monfort (2023) & Psychology and Marketing (3) & Taiwan & 278 food firms (CEO/marketing managers), Taiwan, survey Mar-Aug 2020, RR 33\%; SEM + supplementary fsQCA & survey-SEM & survey construct (implementation of AI marketing strategy) \\
\addlinespace[2pt]
S18 & Cannas et al. (2024) & International Journal of Production Research (3) & Italy & multiple case study: 6 companies, 17 AI implementation cases, semi-structured interviews & case study & case observation (AI applications in SCOR processes) \\
\addlinespace[2pt]
S20 & Pham et al. (2024) & Journal of Business Research (3) & USA & 941 AI hospital-year obs + 941 matched non-AI hospital-year obs (246 matched hospitals), 40 US states, 2000-2020 & panel econometrics & adoption dummy (hospital AI adoption) \\
\addlinespace[2pt]
S21 & Sullivan \& Fosso (2024) & Journal of Business Research (3) & USA & two-wave survey, 107 US executives completing both waves (of 225 wave-1; 47.5\% effective RR), panel-recruited IT+business executives & survey-SEM & survey constructs (AI-enabled automation, AI-enabled analytics, AI-enabled relational capabilities) \\
\addlinespace[2pt]
S24 & Alam et al. (2025) & International Journal of Hospitality Management (3) & Malaysia & 336 respondents, stratified purposive sampling, hospitality industry & survey-SEM & survey construct (AI adoption + technology readiness) \\
\addlinespace[2pt]
S26 & Banna \& Alam (2025) & International Journal of Entrepreneurial Behaviour and Research (3) & EU & unbalanced panel of 1,479 firms (\textasciitilde{}12,900 obs), 26 EU countries, 2012-2023; CGM multi-way clustering + 2SLS-IV & panel econometrics & AI-related venture-capital funding received (log), OECD AI Policy Observatory data (turning point USD 11.3M); robustness: AI job intensity, AI patents, industry-level IV \\
\addlinespace[2pt]
S27 & Basnet et al. (2025) & International Review of Financial Analysis (3) & USA & US firms, 10-K filings 2005-2018 (pre-hype window); AI mentions classified actionable/speculative/irrelevant; causal design & panel econometrics & 10-K text (actionable vs speculative vs irrelevant AI disclosures) \\
\addlinespace[2pt]
S29 & Cao et al. (2025) & European Management Review (3) & USA & Study 1: instrument from 6,519 AI documents of 37 retailers; Study 2: fsQCA on survey of 140 U.S.-based retail managers & fsQCA & validated multi-level instrument (AI integration at task/function/firm level) \\
\addlinespace[2pt]
S32 & D’Amico et al. (2025) & Journal of Technology Transfer (3) & UK & 14,143 UK firms, 2004-2020 & panel econometrics & archival: Beauhurst-identified AI technology/product use, matched to Business Structure Database + UK Innovation Survey (2004-2020, biennial) \\
\addlinespace[2pt]
S34 & Fontanelli et al. (2025) & Journal of Economic Behavior and Organization (3) & France & French ICT survey 2019, firm-level; balancing/matching robustness & panel econometrics & survey (predictive AI use; buyers vs in-house developers) \\
\addlinespace[2pt]
S01 & Wamba-Taguimdje et al. (2020) & Business Process Management Journal (2) & multi-country & \textasciitilde{}500 vendor-published AI project mini-cases (IBM, AWS, Cloudera, Nvidia etc.), archival qualitative analysis & case study & case observation (AI-based transformation projects from vendor websites) \\
\addlinespace[2pt]
S10 & Sun et al. (2022) & Systems Research and Behavioral Science (2) & China & 234 AI-related manufacturers (GM/CEO respondents), CB-SEM & survey-SEM & survey construct (value co-creation in AI innovation ecosystem) \\
\addlinespace[2pt]
S11 & Ali et al. (2023) & Journal of Organizational Change Management (2) & UAE & single case (CMC Dubai): 9 semi-structured interviews with robotic-surgery team + archival data, triangulated & case study & case observation (adoption of AI/robotic surgery) \\
\addlinespace[2pt]
\end{longtable}

\clearpage
\begin{longtable}{@{}l p{2.4cm} l p{4.2cm} l p{6.2cm} p{6.2cm}@{}}
\caption{Included studies: outcomes, effect directions, and conditions (excerpt: 28 of 67 studies --- the 25 in the highest-ranked journals (AJG 4*, 4, 3) plus the first three AJG 2 studies; within a rating by year). Source: author's compilation.}\label{tab:coding-a2}\\
\toprule
ID & Study & Outcome & Outcome measure(s) & Direction & Conditions & Key finding \\
\midrule
\endfirsthead
\multicolumn{7}{@{}l}{\tablename~\thetable{} (continued)}\\
\toprule
ID & Study & Outcome & Outcome measure(s) & Direction & Conditions & Key finding \\
\midrule
\endhead
\bottomrule
\endfoot
S09 & Mishra et al. (2022) & Perf. & Perf: gross and net operating efficiency, net profitability, adspend, employment & mixed &  & AI focus reduces gross operating efficiency but increases net operating efficiency and profitability - automation raises costs more slowly than revenues. \\
\addlinespace[2pt]
S14 & Krakowski et al. (2023) & Both & Perf: competitive performance (game outcomes/ratings); CA: sources of persistent performance heterogeneity (= sustained advantage) across settings & conditional & substitution: traditional human capabilities become obsolete; complementation: new human-machine capabilities emerge that are unrelated or negatively related to traditional ones & AI adoption shifts the sources of competitive advantage - a new decision-making resource emerges at the human-AI intersection, unrelated to traditional capability. \\
\addlinespace[2pt]
S17 & Babina et al. (2024) & Perf. & Perf: growth in sales, employment, and market valuation & positive & channel: product innovation rather than process efficiency; gains concentrate among firms that were already large, which raises industry concentration (superstar dynamics) & AI-investing firms grow in sales, employment and value via product innovation; benefits concentrate among large firms, increasing concentration. \\
\addlinespace[2pt]
S07 & Leoni et al. (2022) & Perf. & Perf: perceived manufacturing firm performance (survey scale) & conditional & full mediation through knowledge management processes: the direct path from AI to manufacturing performance is not significant while both legs of the indirect path are; AI has no direct effect on supply chain resilience either; larger firms reach higher AI maturity & AI alone does not improve manufacturing firm performance - the effect exists exclusively through knowledge management processes (full mediation). \\
\addlinespace[2pt]
S45 & Adams et al. (2026) & Perf. & Perf: growth in sales, employment, assets, markups; stock-return response to monetary policy & positive & adoption concentrated in larger, more productive firms; adopters more responsive to monetary policy surprises & Firms adopting AI-powered pricing grow faster in sales, employment and markups - AI pricing is a capability large productive firms exploit first. \\
\addlinespace[2pt]
S46 & Agag et al. (2026) & Perf. & Perf: firm value; wage inequality as mediating channel & positive & wage inequality partially mediates: AI human capital compresses inequality and thereby raises firm value; the effect is stronger in dynamic and complex industries; by sector, the value effect is strongest in travel and the inequality compression in hospitality & AI human capital raises firm value partly by compressing wage inequality - a workforce-centered pathway to AI-enabled advantage, strongest in dynamic sectors. \\
\addlinespace[2pt]
S52 & Filieri et al. (2026) & Perf. & Perf: supply chain efficiency, sales performance, market value & conditional & inverted U: AI performance benefits diminish at extreme tariff exposure ("optimal investment threshold"); customer engagement needs strengthen the effect (boundary condition); tariff exposure induces AI integration & Trade tensions push firms into AI integration that improves efficiency, sales and value - up to a threshold where adversity overwhelms the compensation. \\
\addlinespace[2pt]
S02 & Baabdullah et al. (2021) & Perf. & Perf: perceived SME performance (survey scale) & positive & technology roadmapping (enabler, sig.); attitude (sig.); infrastructure readiness (sig.); awareness (sig.); professional expertise and technicality not significant & Acceptance of AI practices improves perceived SME performance; adoption is driven by roadmapping, attitude, infrastructure and awareness - not by professional expertise. \\
\addlinespace[2pt]
S03 & Chatterjee et al. (2021) & Both & Perf: perceived organizational performance (survey scale); CA: competitive advantage scale (COA1-3, Rogers-based; incl. market-share gain), discriminant-valid vs OP; predicted by performance & positive & mediators: B2B engagement, employee experience, information processing capability (all sig.); moderator: leadership support (sig., MGA); firm size/age/industry only controls (not moderators) & AI-CRM implementation improves perceived organizational performance and competitive advantage in B2B relationship management. \\
\addlinespace[2pt]
S04 & Chatterjee et al. (2022) & Perf. & Perf: firm performance scale (market+financial criteria, ROI/ROA, sales growth, market share, profitability) + B2B relationship satisfaction & positive & technology turbulence weakens the paths from automated decision-making and operational efficiency to B2B relationship satisfaction (negative moderator); leadership support strengthens the path from satisfaction to firm performance (positive moderator); the main effect of AI-CRM on firm performance is unconditional & AI-embedded CRM improves relationship satisfaction and firm performance; gains shrink under technology turbulence and grow with leadership support. \\
\addlinespace[2pt]
S05 & Hossain et al. (2022) & CA & CA: sustained competitive advantage = relative performance vs competitors over last 3 years (market share growth, sales growth, profitability, ROI) & conditional & mediators: market sensing, seizing and reconfiguring carry the effect of marketing analytics capability on sustained competitive advantage in full; moderator: AI adoption strengthens all three capability paths & Marketing analytics capability drives sensing/seizing/reconfiguring toward sustained competitive advantage; AI adoption amplifies this only when built on an analytics foundation. \\
\addlinespace[2pt]
S06 & Lee et al. (2022) & Perf. & Perf: revenue growth & conditional & AI-intensity threshold (low-level adoption: no effect); complementary technology investments (cloud, database) amplify; internal/exclusive R\&D strategy amplifies & AI pays off only above an adoption-intensity threshold, and more so with complementary technology investments and venture-specific internal R\&D. \\
\addlinespace[2pt]
S08 & Lui et al. (2022) & Perf. & Perf: firm market value (CAR; -1.77\% on event day) & negative & nonmanufacturing firms (more negative); weak IT capability (more negative); low credit rating (more negative) & investors react negatively to AI investment announcements on average; the penalty concentrates in nonmanufacturing firms with weak IT capability and low credit ratings. \\
\addlinespace[2pt]
S12 & Czarnitzki et al. (2023) & Perf. & Perf: firm productivity (labor- and value-added-based) & positive &  & AI adoption is robustly associated with higher firm productivity in representative German data; the effect grows with usage intensity. \\
\addlinespace[2pt]
S15 & Samadhiya et al. (2023) & Both & Perf: perceived firm performance (survey scale); CA: perceived sustainable competitiveness (survey scale) & positive & ICT capability (prerequisite); technological readiness (prerequisite); firm fit & AI-based partner relationship management improves firm performance and sustainable competitiveness, provided ICT capability and technological readiness are in place. \\
\addlinespace[2pt]
S16 & Wu \& Monfort (2023) & Perf. & Perf: perceived firm performance (survey scale) & positive & fsQCA configurations: marketing capabilities + customer value co-creation + market orientation + AI strategy jointly form necessary and sufficient recipes for high performance & AI marketing strategy raises firm performance, but only as part of configurations with marketing capabilities, co-creation and market orientation - a configurational result. \\
\addlinespace[2pt]
S18 & Cannas et al. (2024) & Perf. & Perf: qualitative competitiveness outcomes: costs, lead times, service level, quality, safety, sustainability & positive & barriers as boundary conditions: data quality, AI skills, high investment needs, unclear economic benefits, lack of AI cost-analysis experience & AI improves supply-chain competitiveness across SCOR processes, but realization depends on overcoming data-quality, skill and investment-evaluation barriers. \\
\addlinespace[2pt]
S20 & Pham et al. (2024) & Perf. & Perf: outpatient revenue, inpatient revenue, productivity, occupancy & positive & market share (larger-share hospitals adopt AI and benefit); endogeneity control essential - naive estimates misleading & Hospitals with larger market share adopt AI and gain in revenue, productivity and occupancy; without endogeneity controls the performance effects are misestimated. \\
\addlinespace[2pt]
S21 & Sullivan \& Fosso (2024) & Perf. & Perf: firm performance + process innovation + product innovation (survey scales) & conditional & adaptive response to market change is a full mediator, with no direct AI-to-performance paths modelled; environmental hostility weakens the automation path (negative moderator) and turns the relational path insignificant; environmental dynamism moderates as well; only ten of eighteen conditional indirect effects are significant & AI capabilities pay off through the firms adaptive response to market changes; environment hostility/dynamism condition how much each AI capability contributes. \\
\addlinespace[2pt]
S24 & Alam et al. (2025) & Both & Perf: value creation (survey scale); CA: sustainable competitive advantage (survey scale, modeled as mediator) & positive & sustainable competitive advantage mediates between AI adoption and value creation; dynamic capabilities first impede value creation through transition costs and become crucial later under turbulence; technological turbulence moderates only the path from capabilities to value creation & AI adoption creates value in hospitality mainly through sustainable competitive advantage as transmission channel; dynamic capabilities help only once transition costs are absorbed. \\
\addlinespace[2pt]
S26 & Banna \& Alam (2025) & Perf. & Perf: firm revenue growth & conditional & U-shaped effect with a quantified turning point at roughly USD 11.3M of AI venture funding: below it AI drags revenue, above it AI pays; coupling AI with an R\&D innovation strategy amplifies the effect, standalone AI is insufficient & AI investment follows a U-shaped payoff: short-term revenue drag, long-term gains - and only firms coupling AI with R\&D strategy capture the upside. \\
\addlinespace[2pt]
S27 & Basnet et al. (2025) & Perf. & Perf: firm value (market valuation); R\&D spending, patents, lagged productivity as channels & conditional & what counts is the substance of the disclosure: actionable disclosures bring valuation gains, innovation and lagged productivity, speculative or irrelevant ones bring nothing, and silent or vague peers are penalised & Markets reward only substantive (actionable) AI disclosures - early adopters with real implementation plans gain value via innovation; vague AI talk earns nothing. \\
\addlinespace[2pt]
S29 & Cao et al. (2025) & Perf. & Perf: efficiency and innovation outcomes (configurational) & conditional & configurational: no single AI function suffices - synergistic combinations (esp. customer service + cybersecurity with other functions); emergent capabilities: adaptive learning, predictive analytics, uncertainty mitigation & Retail performance comes from configurations of AI-enabled functions, not isolated AI uses - a directly configurational (fsQCA) result. \\
\addlinespace[2pt]
S32 & D’Amico et al. (2025) & Perf. & Perf: innovation performance (knowledge spillover of innovation) & conditional & complementarity with own R\&D investment (AI pays only if it complements internal R\&D); AI substitutes knowledge collaboration with certain external partners & AI adoption boosts innovation only when paired with the firms own R\&D investment - and partly replaces external knowledge collaboration. \\
\addlinespace[2pt]
S34 & Fontanelli et al. (2025) & Perf. & Perf: volatility of productivity growth (not the level) & conditional & AI buyers: higher volatility; in-house developers: none; complementary human capital (share of ICT engineers/technicians) mitigates the buyer effect & Off-the-shelf AI raises the riskiness of productivity outcomes; developing in-house or holding complementary ICT human capital neutralizes it. \\
\addlinespace[2pt]
S01 & Wamba-Taguimdje et al. (2020) & Perf. & Perf: organizational performance (financial, marketing, administrative) + process-level performance, qualitative & conditional & process reconfiguration required (mechanism: performance gains only when AI features are used to reconfigure processes); AI as configuration of complements (data, talent, domain knowledge, partnerships, infrastructure) & AI transformation projects improve organizational and process performance, but only when firms use AI to reconfigure their processes rather than overlay it on existing ones. \\
\addlinespace[2pt]
S10 & Sun et al. (2022) & CA & CA: perceived competitive advantage scale (own construct, Table 3) + innovation intelligibility; no performance construct in model & conditional & full mediation: value co-creation has no direct effect on competitive advantage; dynamic capabilities and innovation capabilities carry the entire effect & Value co-creation in AI ecosystems has no direct effect on competitive advantage - effects run entirely through dynamic and innovation capabilities. \\
\addlinespace[2pt]
S11 & Ali et al. (2023) & Both & Perf: four qualitative outcome categories: clinical (errors/infections down, patient experience), financial (quantified: \textasciitilde{}500K/month at 30 robotic cases), organizational (resource allocation), technological; CA: competitive position with explicit competitor comparison: patient choice over competitors, bargaining power with third-party payers, differentiation, upper hand over competition (P1) & positive & qualified robotic team incl. revenue-cycle team (capability prerequisite); cultural shift in medical community; regulatory approval (DHA, case-by-case); high acquisition costs not covered by insurance (entry barrier) & AI-enabled robotic surgery improves clinical, financial and technological outcomes and strengthens the competitive position of the healthcare organization. \\
\addlinespace[2pt]
\end{longtable}
\end{landscape}
\endgroup
